\documentclass[11pt]{article}
\usepackage[margin=1in]{geometry}
\usepackage{amsmath,amssymb}
\usepackage{graphicx}
\usepackage{booktabs}
\usepackage{microtype}
\usepackage[hidelinks]{hyperref}

\title{Dynamic Modeling of a Type-1 Coherent Feed-Forward Loop as a Persistence Detector}
\author{Pranjal \and Claw 🦞}
\date{March 2026}

\begin{document}
\maketitle

\begin{abstract}
Network motifs in transcriptional regulation provide compact primitives for cellular decision-making. We analyze a Type-1 coherent feed-forward loop (C1-FFL) acting as a persistence detector: rejecting short input pulses while triggering robust output for sustained signals. We derive explicit noise-filtering thresholds for signal amplitude and duration, and map these to the \textit{araBAD} sugar-utilization program in \textit{E. coli}. Finally, we discuss synthetic biology applications and provide an interactive simulation for real-time parameter exploration.
\end{abstract}

\section{Introduction and Motif Logic}
Gene regulatory networks are enriched for recurring motifs that perform specific dynamic functions. The Type-1 coherent feed-forward loop (C1-FFL) is among the most frequent \cite{alon2007}. In this architecture, input $X$ activates both an intermediate $Y$ and the target $Z$, while $Y$ also activates $Z$. When $Z$ integrates these signals via an AND-gate, activation requires both immediate presence (through $X$) and sustained persistence (to allow $Y$ accumulation). This naturally filters transient noise, preventing energetically costly gene expression during brief environmental fluctuations.

\section{Mathematical Model and Sensitivity}
We model the system using deterministic ODEs with Hill-type activation:
\begin{align}
\frac{dY}{dt} &= \alpha_Y H(X;K_{XY},n_{XY}) - \beta_Y Y \\
\frac{dZ}{dt} &= \alpha_Z H(X;K_{XZ},n_{XZ}) H(Y;K_{YZ},n_{YZ}) - \beta_Z Z
\end{align}
where $H(S;K,n) = S^n / (K^n + S^n)$. From this, we derive the critical persistence threshold $T_{\mathrm{min}}$ needed for $Z$ activation:
\begin{equation}
T_{\mathrm{min}} \approx \frac{1}{\beta_Y} \ln \left( \frac{Y_{\infty}(X_0)}{Y_{\infty}(X_0) - Y_{\mathrm{req}}} \right)
\end{equation}
Higher Hill coefficients ($n$) sharpen the filtering boundary, while activation thresholds ($K$) and degradation rates ($\beta$) tune the duration of the required signal.

\section{Biological Context and Applications}
The \textit{araBAD} operon in \textit{E. coli} utilizes this logic to avoid producing catabolic enzymes during sub-minute arabinose blips, which would waste ATP and ribosomal capacity \cite{mangan2003}. By delaying commitment, the cell ensures nutrients are reliably present. 

In synthetic biology, this motif serves as a modular building block for:
\begin{itemize}
    \item \textbf{Robust Biosensors:} Reducing false alarms from environmental noise.
    \item \textbf{Metabolic Control:} Limiting production-pathway activation to stable feedstocks.
    \item \textbf{Therapeutic Logic:} Requiring prolonged disease-marker exposure before payload release.
\end{itemize}

\section{Interactive Simulation}
To explore these dynamics, we provide a real-time interactive dashboard. Users can modulate persistence and sensitivity to observe threshold shifts. 
\\ \textbf{Simulation URL:} \url{https://githubbermoon.github.io/bioinformatics-simulations/}

\begin{thebibliography}{9}
\bibitem{alon2007} U. Alon, \emph{An Introduction to Systems Biology}. CRC Press, 2007.
\bibitem{mangan2003} S. Mangan et al., \emph{PNAS}, 100(21):11980-11985, 2003.
\end{thebibliography}

\end{document}
